\begin{frame}{Query-Aware Explainable Product Search with Reinforcement Knowledge Graph Reasoning}
\scriptsize
\vspace{-0.8em}

\begin{table}[h!]
\centering
\renewcommand{\arraystretch}{1.5} % increases row height slightly more
\setlength{\tabcolsep}{2.5pt}     % reduces horizontal padding a bit more

\begin{tabular}{|p{2.6cm}|p{4.3cm}|p{4.3cm}|}
\hline
\textbf{Citation} & \textbf{Summary} & \textbf{Advantages \& Disadvantages} \\
\hline
\tiny\textit{[6] Q. Zhu, H. Zhang, Q. He, and Z. Dou, IEEE Trans. Knowl. Data Eng., vol. 36, no. 3, pp. 1260–1273, 2024.  
DOI: 10.1109/TKDE.2023.32973
\newline 31.}
&
Introduces \textbf{QEPS}, a reinforcement learning framework for \textbf{explainable product search} using user-product knowledge graphs. It employs demonstration-guided policy networks and query-aware rewards to generate relevant search results with explanations aligned to user intent. The system also considers user search behaviors to ensure the explanations accurately reflect query-specific reasoning paths.
&
\textbf{Advantages:}
\begin{itemize}\itemsep0.25em
\item Combines \textbf{RL} with knowledge graph reasoning for explainable search.
\item Integrates multiple reward types for better retrieval accuracy.
\item Outperforms baselines on Amazon datasets.
\end{itemize}

\textbf{Disadvantages:}
\begin{itemize}\itemsep0.25em
\item High implementation complexity due to multiple components.
\item Relies on predefined meta-paths, limiting generalizability.
\item \textbf{RL} training is computationally expensive.
\end{itemize} \\
\hline
\end{tabular}
\end{table}
\end{frame}
